\section{Related Work}
\label{sec:related_work}

\paragraph{Self-supervised speech representations for SER.}
WavLM, HuBERT, and emotion-specialized models (e.g., emotion2vec) provide frame-level embeddings that have largely replaced hand-crafted features in SER~\cite{chen2022wavlm,wagner2023dawn}.
Layer-wise analysis shows that mid-to-upper transformer layers often carry the most emotion-discriminative information~\cite{pasad2021layer}, motivating learnable layer fusion rather than fixed layer selection.

\paragraph{Temporal pooling and attention.}
Attention-based pooling aggregates frame sequences into utterance-level vectors but typically operates only in the SSL feature space.
Our prosody-guided variant injects frame-aligned F0 and energy as an explicit prior, related in spirit to knowledge-guided attention in paralinguistics but implemented with strict parameter parity against a self-attention ablation.

\paragraph{Child speech and corpus bias.}
Child SER datasets remain scarce; C-BESD and FAU Aibo represent contrasting regimes (acted vs.\ spontaneous).
Prior work reports strong in-domain acted performance but rarely quantifies cross-style transfer at scale.
We complement parametric Fr\'{e}chet distance (FD) with kernel-based speech MMD (SMMD) to characterize distribution shift non-parametrically.

\paragraph{Explainability in affective computing.}
Post-hoc saliency maps are common, but correlation between learned attention and physical prosody (our APC metric) offers a lightweight, deployment-friendly interpretability probe suitable for clinical review workflows.
